\chapter{2002年京蒙皖卷（春文）}

\section{单选题}

\begin{problem}
不等式组 \(\begin{cases}
x^{2}-1<0,\\
x^{2}-3x<0
\end{cases}\) 的解集为（\quad）

\choices
    {\(\{x\mid -1<x<1\}\)}
    {\(\{x\mid 0<x<3\}\)}
    {\(\{x\mid 0<x<1\}\)}
    {\(\{x\mid -1<x<3\}\)}
\end{problem}

\begin{answer}
C
\end{answer}

\begin{solution}
由 \(x^{2}-1<0\)，得 \(-1<x<1\)；由 \(x^{2}-3x=x\left(x-3\right)<0\)，得 \(0<x<3\). 因此原不等式组的解集为 \(\left(-1,1\right)\cap\left(0,3\right)=\left(0,1\right)\)，故选 C.
\end{solution}

\begin{problem}
已知三条直线 \(m\)，\(n\)，\(l\)，三个平面 \(\alpha\)，\(\beta\)，\(\gamma\)，下列四个命题中，正确的是（\quad）

\choices
    {\(\begin{cases}
\alpha\perp\gamma,\\
\beta\perp\gamma
\end{cases}\Rightarrow\alpha\parallel\beta\)}
    {\(\begin{cases}
m\parallel\beta,\\
l\perp m
\end{cases}\Rightarrow l\perp\beta\)}
    {\(\begin{cases}
m\parallel\gamma,\\
n\parallel\gamma
\end{cases}\Rightarrow m\parallel n\)}
    {\(\begin{cases}
m\perp\gamma,\\
n\perp\gamma
\end{cases}\Rightarrow m\parallel n\)}
\end{problem}

\begin{answer}
D
\end{answer}

\begin{solution}
第一个命题不正确：例如取 \(\gamma\) 为水平面，取两个相交的竖直平面 \(\alpha\)，\(\beta\)，则 \(\alpha\perp\gamma\)，\(\beta\perp\gamma\)，但 \(\alpha\) 与 \(\beta\) 不平行. 第二个命题不正确：直线 \(m\parallel\beta\) 且 \(l\perp m\) 只能说明两条直线的方向垂直，不能保证 \(l\) 垂直于平面 \(\beta\). 第三个命题不正确：两条直线分别平行于同一个平面时，它们的方向可以不同，因而不一定互相平行. 第四个命题正确：垂直于同一个平面的两条直线都与该平面的法线方向平行，所以两直线互相平行，故选 D.
\end{solution}

\begin{problem}
已知椭圆的焦点是 \(F_{1}\)，\(F_{2}\)，\(P\) 是椭圆上的一个动点. 如果延长 \(F_{1}P\) 到 \(Q\)，使得 \(\lvert PQ\rvert=\lvert PF_{2}\rvert\)，那么动点 \(Q\) 的轨迹是（\quad）

\choices
    {圆}
    {椭圆}
    {双曲线的一支}
    {抛物线}
\end{problem}

\begin{answer}
A
\end{answer}

\begin{solution}
设椭圆的长半轴为 \(a\). 因为 \(P\) 在椭圆上，所以 \(F_{1}P+PF_{2}=2a\). 又因为 \(Q\) 在从 \(F_{1}\) 经过 \(P\) 的射线上，且 \(PQ=PF_{2}\)，所以
\[
\lvert F_{1}Q\rvert=\lvert F_{1}P\rvert+\lvert PQ\rvert=\lvert F_{1}P\rvert+\lvert PF_{2}\rvert=2a
\]
因此点 \(Q\) 到定点 \(F_{1}\) 的距离恒为 \(2a\). 又因为焦点 \(F_{1}\) 在椭圆内部，从 \(F_{1}\) 出发的每一条射线都与椭圆相交，故射线 \(F_{1}P\) 的方向可以取遍所有方向，点 \(Q\) 也取遍该圆周. 所以 \(Q\) 的轨迹是以 \(F_{1}\) 为圆心、半径 \(2a\) 的圆，故选 A.
\end{solution}

\begin{problem}
如果 \(\theta\in\left(\dfrac{\pi}{2},\pi\right)\)，那么复数 \(\left(1+\i\right)\left(\cos\theta+\i\sin\theta\right)\) 的辐角的主值是（\quad）

\choices
    {\(\theta+\dfrac{9\pi}{4}\)}
    {\(\theta+\dfrac{\pi}{4}\)}
    {\(\theta-\dfrac{\pi}{4}\)}
    {\(\theta+\dfrac{7\pi}{4}\)}
\end{problem}

\begin{answer}
B
\end{answer}

\begin{solution}
有 \(1+\i=\sqrt{2}\left(\cos\dfrac{\pi}{4}+\i\sin\dfrac{\pi}{4}\right)\)，所以
\(\left(1+\i\right)\left(\cos\theta+\i\sin\theta\right)=\sqrt{2}\left(\cos\left(\theta+\dfrac{\pi}{4}\right)+\i\sin\left(\theta+\dfrac{\pi}{4}\right)\right)\).
又因为 \(\dfrac{\pi}{2}<\theta<\pi\)，所以 \(\dfrac{3\pi}{4}<\theta+\dfrac{\pi}{4}<\dfrac{5\pi}{4}\)，该角在 \([0,2\pi)\) 内，故辐角的主值为 \(\theta+\dfrac{\pi}{4}\)，选 B.
\end{solution}

\begin{problem}
若角 \(\alpha\) 满足条件 \(\sin2\alpha<0\)，\(\cos\alpha-\sin\alpha<0\)，则 \(\alpha\) 在（\quad）

\choices
    {第一象限}
    {第二象限}
    {第三象限}
    {第四象限}
\end{problem}

\begin{answer}
B
\end{answer}

\begin{solution}
在一个周期内，\(\sin2\alpha<0\) 给出 \(\alpha\) 位于第二或第四象限. 第二象限内 \(\cos\alpha<0\)，\(\sin\alpha>0\)，从而 \(\cos\alpha-\sin\alpha<0\) 成立；第四象限内 \(\cos\alpha>0\)，\(\sin\alpha<0\)，从而 \(\cos\alpha-\sin\alpha>0\)，不满足题设. 因此 \(\alpha\) 在第二象限，故选 B.
\end{solution}

\begin{problem}
从 6 名志愿者中选出 4 人分别从事翻译，导游，导购，保洁四项不同工作. 若其中甲，乙两名志愿者都不能从事翻译工作，则选派方案共有（\quad）

\choices
    {\(280\text{ 种}\)}
    {\(240\text{ 种}\)}
    {\(180\text{ 种}\)}
    {\(96\text{ 种}\)}
\end{problem}

\begin{answer}
B
\end{answer}

\begin{solution}
翻译工作不能由甲、乙承担，所以翻译人员有 \(4\) 种选法. 翻译人员确定后，从剩下的 \(5\) 人中选出 \(3\) 人分别担任导游、导购、保洁，有 \(\mathrm A_{5}^{3}=5\times4\times3=60\) 种选法. 因此选派方案共有 \(4\times60=240\) 种，故选 B.
\end{solution}

\begin{problem}
在 \(\triangle ABC\) 中，\(AB=2\)，\(BC=1.5\)，\(\angle ABC=120^{\circ}\). 若使三角形绕直线 \(BC\) 旋转一周，则所形成的几何体的体积是（\quad）

\choices
    {\(\dfrac{3}{2}\pi\)}
    {\(\dfrac{5}{2}\pi\)}
    {\(\dfrac{7}{2}\pi\)}
    {\(\dfrac{9}{2}\pi\)}
\end{problem}

\begin{answer}
A
\end{answer}

\begin{solution}
设从 \(A\) 向直线 \(BC\) 作垂线，垂足为 \(H\). 因为 \(AB=2\)，\(\angle ABC=120^{\circ}\)，所以 \(AH=AB\sin120^{\circ}=\sqrt{3}\)，又 \(BH=\lvert AB\cos120^{\circ}\rvert=1\).
点 \(H\) 在 \(B\) 的外侧，故 \(CH=BC+BH=\dfrac{5}{2}\). 绕 \(BC\) 旋转后，可看成以 \(C\) 为顶点、底面半径为 \(AH\)、高为 \(CH\) 的圆锥，减去以 \(B\) 为顶点、底面半径为 \(AH\)、高为 \(BH\) 的圆锥. 因此
\[
V=\dfrac{1}{3}\pi AH^{2}(CH-BH)=\dfrac{1}{3}\pi\cdot3\cdot\dfrac{3}{2}=\dfrac{3}{2}\pi
\]
故选 A.
\end{solution}

\begin{problem}
到两坐标轴距离相等的点的轨迹方程是（\quad）

\choices
    {\(x-y=0\)}
    {\(x+y=0\)}
    {\(\lvert x\rvert-y=0\)}
    {\(\lvert x\rvert-\lvert y\rvert=0\)}
\end{problem}

\begin{answer}
D
\end{answer}

\begin{solution}
点 \((x,y)\) 到 \(x\) 轴的距离为 \(\lvert y\rvert\)，到 \(y\) 轴的距离为 \(\lvert x\rvert\). 两距离相等的条件是 \(\lvert x\rvert=\lvert y\rvert\)，即 \(\lvert x\rvert-\lvert y\rvert=0\)，故选 D.
\end{solution}

\begin{problem}
函数 \(y=2^{\sin x}\) 的单调区间是（\quad）

\choices
    {\(\left[2k\pi-\dfrac{\pi}{2},2k\pi+\dfrac{\pi}{2}\right]\)（\(k\in\Z\)）}
    {\(\left[2k\pi+\dfrac{\pi}{2},2k\pi+\dfrac{3\pi}{2}\right]\)（\(k\in\Z\)）}
    {\([2k\pi-\pi,2k\pi]\)（\(k\in\Z\)）}
    {\([2k\pi,2k\pi+\pi]\)（\(k\in\Z\)）}
\end{problem}

\begin{answer}
A
\end{answer}

\begin{solution}
因为底数 \(2>1\)，函数 \(y=2^{u}\) 随 \(u\) 增大而增大，所以 \(y=2^{\sin x}\) 与 \(y=\sin x\) 具有相同的单调性. 正弦函数的单调递增区间为
\(\left[2k\pi-\dfrac{\pi}{2},2k\pi+\dfrac{\pi}{2}\right]\)（\(k\in\Z\)），正弦函数的单调递减区间为
\(\left[2k\pi+\dfrac{\pi}{2},2k\pi+\dfrac{3\pi}{2}\right]\)（\(k\in\Z\)）.
原始试卷的题干写作“单调区间”. 按此表述，第一项是单调递增区间，第二项是单调递减区间，二者均为单调区间，题目不能唯一确定选项. 公开同题资料将题干写作“单调增区间”，若按这一校正版题意理解，则应选第一项；现稿答案依原卷公开答案保留，并记录该题面歧义.
\end{solution}

\begin{problem}
在 \(\left(\dfrac{1}{x}+x^{2}\right)^{6}\) 的展开式中，\(x^{3}\) 的系数和常数项依次是（\quad）

\choices
    {\(20\)，\(20\)}
    {\(15\)，\(20\)}
    {\(20\)，\(15\)}
    {\(15\)，\(15\)}
\end{problem}

\begin{answer}
C
\end{answer}

\begin{solution}
展开式的通项为
\(\mathrm C_{6}^{k}\left(\dfrac{1}{x}\right)^{6-k}(x^{2})^{k}=\mathrm C_{6}^{k}x^{3k-6}\)，\(k=0\)，\(1\)，\(\ldots\)，\(6\)
当幂指数为 \(3\) 时，\(3k-6=3\)，得 \(k=3\)，所以 \(x^{3}\) 的系数为 \(\mathrm C_{6}^{3}=20\). 当幂指数为 \(0\) 时，\(3k-6=0\)，得 \(k=2\)，所以常数项为 \(\mathrm C_{6}^{2}=15\). 故选 C.
\end{solution}

\begin{problem}
若一个等差数列前 3 项的和为 34，最后 3 项的和为 146，且所有项的和为 390，则这个数列有（\quad）

\choices
    {\(13\text{ 项}\)}
    {\(12\text{ 项}\)}
    {\(11\text{ 项}\)}
    {\(10\text{ 项}\)}
\end{problem}

\begin{answer}
A
\end{answer}

\begin{solution}
设等差数列的公差为 \(d\)，中间项分别为 \(a_{2}\) 和 \(a_{n-1}\). 由等差数列性质，前 \(3\) 项的和为 \(3a_{2}\)，最后 \(3\) 项的和为 \(3a_{n-1}\)，所以
\(a_{2}=\dfrac{34}{3}\)，又 \(a_{n-1}=\dfrac{146}{3}\). 又有 \(a_{1}+a_{n}=a_{2}+a_{n-1}=60\)，于是数列的前 \(n\) 项和为 \(S_{n}=\dfrac{n}{2}\left(a_{1}+a_{n}\right)=30n\).
由 \(S_{n}=390\)，得 \(30n=390\)，即 \(n=13\)，故选 A.
\end{solution}

\begin{problem}
用一张钢板制作一个容积为 \(4\mathrm{~m}^{3}\) 的无盖长方体水箱. 可用的长方形钢板有四种不同的规格（长\(\times\)宽的尺寸如各选项所示，单位均为 \(\mathrm{m}\)），若既要够用，又要所剩最少，则应选钢板的规格是（\quad）

\choices
    {\(2\times5\)}
    {\(2\times5.5\)}
    {\(2\times6.1\)}
    {\(3\times5\)}
\end{problem}

\begin{answer}
C
\end{answer}

\begin{solution}
设水箱底面长、宽和高分别为 \(x\)，\(y\)，\(z\)，则 \(xyz=4\). 无盖水箱所需钢板面积为
\[
S=xy+2xz+2yz
\]
由基本不等式，
\[
S\geqslant3\sqrt[3]{\left(xy\right)\left(2xz\right)\left(2yz\right)}=3\sqrt[3]{4x^{2}y^{2}z^{2}}=3\sqrt[3]{64}=12
\]
当 \(xy=2xz=2yz\) 时取等号，从而 \(x=y=2z\)，结合 \(xyz=4\) 得 \(x=y=2\)，\(z=1\). 因此制作容积为 \(4\mathrm{~m}^{3}\) 的无盖水箱至少需要 \(12\mathrm{~m}^{2}\) 钢板.

四种钢板面积分别为 \(10\mathrm{~m}^{2}\)，\(11\mathrm{~m}^{2}\)，\(12.2\mathrm{~m}^{2}\)，\(15\mathrm{~m}^{2}\). 规格 \(2\times6.1\) 可以沿长边排下一个 \(2\times2\) 的底板和四个 \(2\times1\) 的侧板，确实够用. 规格 \(3\times5\) 可在宽为 \(2\) 的区域内沿长边排下一个 \(2\times2\) 的底板和三个 \(2\times1\) 的侧板，再将第四个侧板转置放入剩余的 \(1\times5\) 区域内，也确实够用. 因此只有规格 \(2\times6.1\) 和 \(3\times5\) 够用，所剩面积分别为 \(0.2\mathrm{~m}^{2}\) 和 \(3\mathrm{~m}^{2}\)，所以应选规格 \(2\times6.1\)，故选 C.
\end{solution}

\section{填空题}

\begin{problem}
若双曲线 \(\dfrac{x^{2}}{4}-\dfrac{y^{2}}{m}=1\) 的渐近线方程为 \(y=\pm\dfrac{\sqrt{3}}{2}x\)，则双曲线的焦点坐标是\fillinblank{}.
\end{problem}

\begin{answer}
\(\left(-\sqrt{7},0\right)\)，\(\left(\sqrt{7},0\right)\)
\end{answer}

\begin{solution}
双曲线的标准形式中，\(a^{2}=4\)，所以 \(a=2\)，渐近线斜率的绝对值为 \(\dfrac{b}{a}\). 由
\(\dfrac{b}{2}=\dfrac{\sqrt{3}}{2}\)，得 \(b=\sqrt{3}\)，从而 \(m=b^{2}=3\). 又因为 \(c^{2}=a^{2}+b^{2}=4+3=7\)，所以两个焦点坐标为 \(\left(-\sqrt{7},0\right)\)，\(\left(\sqrt{7},0\right)\).
\end{solution}

\begin{problem}
如果 \(\cos\theta=-\dfrac{12}{13}\)，\(\theta\in\left(\pi,\dfrac{3\pi}{2}\right)\)，那么 \(\cos\left(\theta+\dfrac{\pi}{4}\right)\) 的值等于\fillinblank{}.
\end{problem}

\begin{answer}
\(-\dfrac{7\sqrt{2}}{26}\)
\end{answer}

\begin{solution}
因为 \(\theta\) 在第三象限，且 \(\cos\theta=-\dfrac{12}{13}\)，所以 \(\sin\theta=-\sqrt{1-\cos^{2}\theta}=-\sqrt{1-\dfrac{144}{169}}=-\dfrac{5}{13}\).
由和角公式，
\[
\cos\left(\theta+\dfrac{\pi}{4}\right)=\cos\theta\cos\dfrac{\pi}{4}-\sin\theta\sin\dfrac{\pi}{4}=\left(-\dfrac{12}{13}+\dfrac{5}{13}\right)\dfrac{\sqrt{2}}{2}=-\dfrac{7\sqrt{2}}{26}
\]
\end{solution}

\begin{problem}
正方形 \(ABCD\) 的边长是 2，\(E\)，\(F\) 分别是 \(AB\) 和 \(CD\) 的中点，将正方形沿 \(EF\) 折成直二面角（如图所示）. \(M\) 为矩形 \(AEFD\) 内一点，如果 \(\angle MBE=\angle MBC\)，\(MB\) 和平面 \(BCF\) 所成角的正切值为 \(\dfrac{1}{2}\)，那么点 \(M\) 到直线 \(EF\) 的距离为\fillinblank{}.

\bitmapfigure[width=0.48\linewidth]{img/2002/jing_meng_wan_liberal_spring/q15_fig1.png}
\end{problem}

\begin{answer}
\(\dfrac{\sqrt{2}}{2}\)
\end{answer}

\begin{solution}
折叠后取 \(EF\) 所在直线为 \(x\) 轴，建立空间直角坐标系，使 \(E=(0,0,0)\)，\(F=(2,0,0)\)，\(B=(0,1,0)\)，\(C=(2,1,0)\)，且平面 \(BCF\) 为 \(z=0\). 由于 \(AE=EB=1\)，矩形 \(AEFD\) 所在平面与平面 \(BCF\) 垂直，设 \(M=(t,0,h)\)，其中 \(0<t<2\)，\(0<h<1\).

于是 \(\overrightarrow{BM}=(t,-1,h)\)，\(\overrightarrow{BE}=(0,-1,0)\)，\(\overrightarrow{BC}=(2,0,0)\). 由 \(\angle MBE=\angle MBC\)，有
\[
\cos\angle MBE=\dfrac{\overrightarrow{BM}\cdot\overrightarrow{BE}}{\lvert BM\rvert\,\lvert BE\rvert}=\dfrac{1}{\sqrt{t^{2}+1+h^{2}}}
\]
以及
\[
\cos\angle MBC=\dfrac{\overrightarrow{BM}\cdot\overrightarrow{BC}}{\lvert BM\rvert\,\lvert BC\rvert}=\dfrac{t}{\sqrt{t^{2}+1+h^{2}}}
\]
因此 \(t=1\)，即 \(M\) 在平面 \(AEFD\) 内位于过 \(F\) 且垂直于 \(EF\) 的直线上. 直线 \(MB\) 在平面 \(BCF\) 内的射影长度为 \(\sqrt{t^{2}+1}=\sqrt{2}\).
设点 \(M\) 到平面 \(BCF\) 的距离为 \(h\)，则直线 \(MB\) 与平面 \(BCF\) 所成角的正切值为
 \(\tan\angle\left(MB,\text{平面 }BCF\right)=\dfrac{h}{\sqrt{2}}=\dfrac{1}{2}\).
所以 \(h=\dfrac{\sqrt{2}}{2}\). 又因为 \(M=(1,0,h)\)，所以点 \(M\) 到直线 \(EF\) 的距离为 \(h=\dfrac{\sqrt{2}}{2}\).
\end{solution}

\begin{problem}
对于任意两个复数 \(z_{1}=x_{1}+y_{1}\i\)，\(z_{2}=x_{2}+y_{2}\i\)（\(x_{1}\)，\(y_{1}\)，\(x_{2}\)，\(y_{2}\) 为实数），定义运算 \(\odot\) 为：\(z_{1}\odot z_{2}=x_{1}x_{2}+y_{1}y_{2}\). 设非零复数 \(w_{1}\)，\(w_{2}\) 在复平面内对应的点分别为 \(P_{1}\)，\(P_{2}\)，点 \(O\) 为坐标原点. 如果 \(w_{1}\odot w_{2}=0\)，那么在 \(\triangle P_{1}OP_{2}\) 中，\(\angle P_{1}OP_{2}\) 的大小为\fillinblank{}.
\end{problem}

\begin{answer}
\(90^{\circ}\)
\end{answer}

\begin{solution}
非零复数 \(w_{1}\)，\(w_{2}\) 对应的向量分别为 \(\overrightarrow{OP_{1}}=(x_{1},y_{1})\)，\(\overrightarrow{OP_{2}}=(x_{2},y_{2})\). 题设运算条件给出
 \(\overrightarrow{OP_{1}}\cdot\overrightarrow{OP_{2}}=x_{1}x_{2}+y_{1}y_{2}=w_{1}\odot w_{2}=0\).
由于两个向量都非零，它们的内积为零等价于两向量互相垂直. 因此 \(\angle P_{1}OP_{2}=90^{\circ}\).
\end{solution}

\section{解答题}

\begin{problem}
在 \(\triangle ABC\) 中，已知 \(A\)，\(B\)，\(C\) 成等差数列，求 \(\tan\dfrac{A}{2}+\tan\dfrac{C}{2}+\sqrt{3}\tan\dfrac{A}{2}\tan\dfrac{C}{2}\) 的值.
\end{problem}

\begin{solution}
因为 \(A\)，\(B\)，\(C\) 成等差数列，所以 \(A+C=2B\). 又因为 \(A+B+C=180^{\circ}\)，所以 \(3B=180^{\circ}\)，即 \(B=60^{\circ}\)，从而
 \(\dfrac{A}{2}+\dfrac{C}{2}=\dfrac{A+C}{2}=B=60^{\circ}\).
设 \(u=\tan\dfrac{A}{2}\)，\(v=\tan\dfrac{C}{2}\). 由正切的和角公式，
\[
\sqrt{3}=\tan60^{\circ}=\tan\left(\dfrac{A}{2}+\dfrac{C}{2}\right)=\dfrac{u+v}{1-uv}
\]
由于 \(0<\dfrac{A}{2}+\dfrac{C}{2}<90^{\circ}\)，分母不为零，所以 \(u+v=\sqrt{3}(1-uv)\). 因此
\[
\tan\dfrac{A}{2}+\tan\dfrac{C}{2}+\sqrt{3}\tan\dfrac{A}{2}\tan\dfrac{C}{2}=u+v+\sqrt{3}uv=\sqrt{3}
\]
\end{solution}

\begin{problem}
已知函数 \(f(x)\) 是偶函数，而且在 \((0,+\infty)\) 上是增函数，判断 \(f(x)\) 在 \((-\infty,0)\) 上是增函数还是减函数，并证明你的判断.
\end{problem}

\begin{solution}
设 \(x_{1}<x_{2}<0\). 则 \(-x_{1}>-x_{2}>0\). 因为 \(f(x)\) 在 \((0,+\infty)\) 上是增函数，所以 \(f(-x_{1})>f(-x_{2})\). 又因为 \(f(x)\) 是偶函数，\(f(x_{1})=f(-x_{1})\)，\(f(x_{2})=f(-x_{2})\)，于是 \(f(x_{1})>f(x_{2})\). 因此当 \(x_{1}<x_{2}<0\) 时有 \(f(x_{1})>f(x_{2})\)，故 \(f(x)\) 在 \((-\infty,0)\) 上是减函数.
\end{solution}

\begin{problem}
在三棱锥 \(S-ABC\) 中，如图，\(\angle SAB=\angle SAC=\angle ACB=90^{\circ}\)，\(AC=2\)，\(BC=\sqrt{13}\)，\(SB=\sqrt{29}\).

\begin{enumerate}
\item 证明：\(SC\perp BC\)；
\item 求侧面 \(SBC\) 与底面 \(ABC\) 所成的二面角大小；
\item 求三棱锥的体积 \(V_{S-ABC}\).
\end{enumerate}

\bitmapfigure[width=0.40\linewidth]{img/2002/jing_meng_wan_liberal_spring/q19_fig1.png}
\end{problem}

\begin{solution}
\begin{enumerate}
\item 因为 \(\angle SAB=\angle SAC=90^{\circ}\)，且 \(AB\)，\(AC\) 是平面 \(ABC\) 内相交的两条直线，所以 \(SA\perp\) 平面 \(ABC\). 又因为 \(\angle ACB=90^{\circ}\)，所以
 \(AB^{2}=AC^{2}+BC^{2}=2^{2}+\left(\sqrt{13}\right)^{2}=17\). 在直角三角形 \(SAB\) 中，\(SA^{2}=SB^{2}-AB^{2}=29-17=12\). 在直角三角形 \(SAC\) 中，\(SC^{2}=SA^{2}+AC^{2}=12+4=16\).
于是 \(SC^{2}+BC^{2}=16+13=29=SB^{2}\). 根据勾股定理的逆定理，\(\angle SCB=90^{\circ}\)，所以 \(SC\perp BC\).

 \item 在平面 \(SBC\) 内，\(SC\perp BC\)；在平面 \(ABC\) 内，\(AC\perp BC\). 因此侧面 \(SBC\) 与底面 \(ABC\) 所成的二面角等于 \(\angle SCA\). 在直角三角形 \(SAC\) 中，\(\tan\angle SCA=\dfrac{SA}{AC}=\dfrac{2\sqrt{3}}{2}=\sqrt{3}\).
所以二面角的大小为 \(60^{\circ}\).

 \item 底面 \(ABC\) 是直角三角形，且 \(SA\perp\) 平面 \(ABC\)，所以 \(SA\) 是三棱锥的高. 底面面积为 \(S_{\triangle ABC}=\dfrac{1}{2}AC\cdot BC=\dfrac{1}{2}\cdot2\cdot\sqrt{13}=\sqrt{13}\).
故
\[
V_{S-ABC}=\dfrac{1}{3}S_{\triangle ABC}\cdot SA=\dfrac{1}{3}\cdot\sqrt{13}\cdot2\sqrt{3}=\dfrac{2\sqrt{39}}{3}
\]
\end{enumerate}
\end{solution}

\begin{problem}
假设 \(A\) 型进口车关税税率在 \(2001\) 年是 \(100\%\)，在 \(2006\) 年是 \(25\%\)，\(2001\) 年 \(A\) 型进口车每辆价格为 \(64\text{ 万元}\)（其中含 \(32\text{ 万元}\)关税税款）.

\begin{enumerate}
\item 已知与 \(A\) 型车性能相近的 \(B\) 型国产车，\(2001\) 年每辆价格为 \(46\text{ 万元}\)，若 \(A\) 型车的价格只受关税降低的影响，为了保证 \(2006\) 年 \(B\) 型车的价格不高于 \(A\) 型车价格的 90\%，\(B\) 型车价格要逐年降低，问平均每年至少下降多少万元？
\item 某人在 \(2001\) 年将 \(33\text{ 万元}\)存入银行，假设银行扣利息税后的年利率为 \(1.8\%\)（\(5\) 年内不变），且每年按复利计算（上一年的利息计入第二年的本金），那么五年到期时这笔钱连本带息是否一定够买按（1）中所述降价后的 \(B\) 型车一辆？
\end{enumerate}
\end{problem}

\begin{solution}
\begin{enumerate}
\item \(2001\) 年 \(A\) 型车的去税价格为 \(64-32=32\) 万元，因此 \(2006\) 年关税税率为 \(25\%\) 时，\(A\) 型车价格为 \(32\left(1+25\%\right)=40\) 万元. 其价格的 \(90\%\) 为 \(36\) 万元. 设 \(B\) 型车平均每年下降 \(x\) 万元，则五年后的价格不高于 \(46-5x\)，要满足题意需有
\[
46-5x\leqslant36
\]
即 \(x\geqslant2\). 所以平均每年至少下降 \(2\) 万元.

\item 按第（1）问的最低降价标准，\(2006\) 年 \(B\) 型车价格不超过 \(36\) 万元. 存款五年后的本息为
\[
33\left(1+1.8\%\right)^{5}=33\left(1.018\right)^{5}\approx36.0789>36
\]
因此无论 \(B\) 型车实际降价幅度是否超过最低标准，这笔钱五年到期时都一定够买一辆降价后的 \(B\) 型车.
\end{enumerate}
\end{solution}

\begin{problem}
已知某椭圆的焦点是 \(F_{1}(-4,0)\)，\(F_{2}(4,0)\)，过点 \(F_{2}\)，并垂直于 \(x\) 轴的直线与椭圆的一个交点为 \(B\)，且 \(\lvert F_{1}B\rvert+\lvert F_{2}B\rvert=10\). 椭圆上不同的两点 \(A(x_{1},y_{1})\)，\(C(x_{2},y_{2})\) 满足条件：\(\lvert F_{2}A\rvert\)，\(\lvert F_{2}B\rvert\)，\(\lvert F_{2}C\rvert\) 成等差数列.

\begin{enumerate}
\item 求该椭圆的方程；
\item 求弦 \(AC\) 中点的横坐标.
\end{enumerate}
\end{problem}

\begin{solution}
\begin{enumerate}
\item 焦点为 \(F_{1}(-4,0)\)，\(F_{2}(4,0)\)，所以半焦距 \(c=4\). 椭圆上任意一点到两焦点的距离之和为 \(2a\)，由点 \(B\) 的条件得 \(2a=10\)，即 \(a=5\). 因而
\[
b^{2}=a^{2}-c^{2}=25-16=9
\]
所以椭圆方程为
\[
\dfrac{x^{2}}{25}+\dfrac{y^{2}}{9}=1
\]

\item 直线 \(x=4\) 与椭圆相交于 \(B\)，代入方程得
\[
\dfrac{16}{25}+\dfrac{y_{B}^{2}}{9}=1
\]
所以 \(y_{B}^{2}=\dfrac{81}{25}\).
所以 \(\lvert F_{2}B\rvert=\lvert y_{B}\rvert=\dfrac{9}{5}\). 记 \(r_{A}=\lvert F_{2}A\rvert\)，\(r_{C}=\lvert F_{2}C\rvert\)，由等差数列条件，
\[
r_{A}+r_{C}=2\lvert F_{2}B\rvert=\dfrac{18}{5}
\]
对椭圆上任意点 \(P(x,y)\)，记 \(r=\lvert F_{2}P\rvert\). 由椭圆定义，\(\lvert F_{1}P\rvert=10-r\)，于是
\[
\left(10-r\right)^{2}-r^{2}=\lvert F_{1}P\rvert^{2}-\lvert F_{2}P\rvert^{2}=16x
\]
化简得 \(100-20r=16x\)，即 \(x=\dfrac{25-5r}{4}\). 因此
\[
x_{1}+x_{2}=\dfrac{25-5r_{A}}{4}+\dfrac{25-5r_{C}}{4}=\dfrac{50-5\left(r_{A}+r_{C}\right)}{4}=8
\]
所以弦 \(AC\) 中点的横坐标为 \(\dfrac{x_{1}+x_{2}}{2}=4\).
\end{enumerate}
\end{solution}

\begin{problem}
已知点的序列 \(A_{n}(x_{n},0)\)，\(n\in\N\)，其中 \(x_{1}=0\)，\(x_{2}=a\)（\(a<0\)），\(A_{3}\) 是线段 \(A_{1}A_{2}\) 的中点，\(A_{4}\) 是线段 \(A_{2}A_{3}\) 的中点，\(\cdots\)，\(A_{n}\) 是线段 \(A_{n-2}A_{n-1}\) 的中点，\(\cdots\).

\begin{enumerate}
\item 写出 \(x_{n}\) 与 \(x_{n-1}\)，\(x_{n-2}\) 之间的关系式（\(n\geqslant3\)）；
\item 设 \(a_{n}=x_{n+1}-x_{n}\)，计算 \(a_{1}\)，\(a_{2}\)，\(a_{3}\)，由此推测数列 \(\{a_{n}\}\) 的通项公式，并加以证明；
\item 求 \(\displaystyle\lim\limits_{n\to\infty}x_{n}\).
\end{enumerate}
\end{problem}

\begin{solution}
\begin{enumerate}
\item 由中点定义，\(A_{n}\) 是 \(A_{n-2}A_{n-1}\) 的中点，所以
 \(x_{n}=\dfrac{x_{n-2}+x_{n-1}}{2}\)，其中 \(n\geqslant3\).

\item 由 \(x_{1}=0\)，\(x_{2}=a\)，得 \(a_{1}=x_{2}-x_{1}=a\). 又 \(x_{3}=\dfrac{x_{1}+x_{2}}{2}=\dfrac{a}{2}\)，所以 \(a_{2}=x_{3}-x_{2}=-\dfrac{a}{2}\). 再由 \(x_{4}=\dfrac{x_{2}+x_{3}}{2}=\dfrac{3a}{4}\)，得 \(a_{3}=x_{4}-x_{3}=\dfrac{a}{4}\).
由第（1）问的递推关系，
\[
a_{n}=x_{n+1}-x_{n}=\dfrac{x_{n-1}+x_{n}}{2}-x_{n}=-\dfrac{1}{2}(x_{n}-x_{n-1})=-\dfrac{1}{2}a_{n-1}
\]
因此由数学归纳法，
\(a_{n}=\dfrac{(-1)^{n+1}a}{2^{n-1}}\)，\(n\geqslant1\)
事实上，\(n=1\) 时公式给出 \(a_{1}=a\)；若对某个 \(n\) 成立，则 \(a_{n+1}=-\dfrac12a_{n}=\dfrac{(-1)^{n+2}a}{2^{n}}\)，故公式对所有正整数 \(n\) 成立.

\item 因为 \(x_{n}=x_{1}+\sum_{k=1}^{n-1}a_{k}\)，所以
\[
x_{n}=a\sum_{k=0}^{n-2}\left(-\dfrac12\right)^{k}=\dfrac{2a}{3}\left(1-\left(-\dfrac12\right)^{n-1}\right)
\]
由于 \(\left(-\dfrac12\right)^{n-1}\to0\)，故
\[
\displaystyle\lim\limits_{n\to\infty}x_{n}=\dfrac{2a}{3}
\]
\end{enumerate}
\end{solution}
